\documentclass[a4paper,titlepage,halfparskip,12pt]{scrreprt}

\usepackage[ngerman]{babel, varioref}
\usepackage[utf8]{inputenc}
\usepackage[T1]{fontenc}
\usepackage{graphicx}
\usepackage{fancyhdr}
\usepackage{amsmath}
\usepackage{geometry}
\geometry{a4paper, top=25mm,left=25mm,right=25mm,bottom=25mm, footskip=12mm}
\usepackage{longtable}
\usepackage{setspace}
\usepackage{lmodern}
%blocksatz
\sloppy
%formatierung literaturverzeichnisangabe
\bibliographystyle{unsrt}

%Auflistungen von Punkten
\usepackage{paralist} 
%urls anzeigen
\usepackage{url}

%Codelisting
\usepackage{xcolor}
\definecolor{mygreen}{rgb}{0,0.6,0}
\definecolor{mygray}{rgb}{0.5,0.5,0.5}
\definecolor{mymauve}{rgb}{0.58,0,0.82}
\definecolor{burntorange}{rgb}{0.8, 0.33, 0.0}
\definecolor{cornellred}{rgb}{0.7, 0.11, 0.11}

\usepackage{listingsutf8}
\lstset{
commentstyle=\color{mygreen},
numberstyle=\small\color{black},
stringstyle=\color{mymauve},
emph={square}, 
showstringspaces=false,
flexiblecolumns=false,
tabsize=2,
numbers=left,
numberblanklines=false,
stepnumber=1,
captionpos=b,
numbersep=5pt,
xleftmargin=15pt,
breaklines=true,
inputencoding=utf8,
extendedchars=true,
extendedchars=true,
basicstyle=\ttfamily\footnotesize,
keywordstyle = \bfseries\color{burntorange},
keywordstyle = [2]\bfseries\color{cornellred},
literate=%
    {Ä}{{\"A}}1%
    {Ö}{{\"O}}1%
    {Ü}{{\"U}}1%
    {ä}{{\"a}}1%
    {ö}{{\"o}}1%
    {ü}{{\"u}}1%
    {ß}{{\ss}}1,%
frame=single,
frameround=ffff
}

%meta data
\usepackage[hidelinks]{hyperref}
\urlstyle{same}

%akronymverzeichnis
\usepackage[printonlyused]{acronym}

% titel definieren
\newcommand{\titel}{Entwicklung eines Chatsystems\\auf Basis von XMPP}

%autor definieren
\newcommand{\autor}{Lukas Priester,Oliver Klapper}
\newcommand{\keywords}{\autor,\titel,Studienarbeit}

% Allgemeines für das PDF
\hypersetup{
    pdftitle={\titel},
    pdfauthor={\autor},
    pdfcreator={\autor},
    pdfsubject={\titel},
    pdflang={Deutsch},
    pdfdisplaydoctitle=true,
    pdfkeywords={\keywords},
}

% set distances of chapter headlines in document
\renewcommand*\chapterheadstartvskip{\vspace*{20pt}} % set distance to header
% set distance to text
%\renewcommand*\chapterheadendvskip{%
%  \vspace*{1\baselineskip plus .1\baselineskip minus .167\baselineskip}}


\begin{document}

\begin{table}[h]
\centering
\begin{tabular}{lcr}
\includegraphics[height=3.5cm]{images/dhbw-logo}
\end{tabular}
\end{table}
\bigskip
\bigskip
\begin{center}
\vspace*{12mm} {\LARGE\textbf{\titel}}\\
\vspace*{12mm} {\large\textbf{Studienarbeit}}\\
\vspace*{3mm} {\large\textbf{5. - 6. Semester}}\\
\vspace*{12mm} des Studiengangs Informationstechnik (B.Sc.)\\ an der Dualen Hochschule Baden-Württemberg Stuttgart\\
% \vspace*{3mm} an der Dualen Hochschule Baden-Württemberg\\
\vspace*{12mm} von\\
\vspace*{3mm} {\large\textbf{Lukas Priester, Oliver Klapper}}\\
\vspace*{12mm} \today\\
\end{center}
\vfill
\begin{spacing}{1.5}
\begin{tabbing}
mmmmmmmmmmmmmmmmmmmmmmmmmm \= \kill
\textbf{Bearbeitungszeitraum} \> 01.10.2019 - 01.05.2020\\
\textbf{Matrikelnummer, Kurs} \> 7288057, 4191693 \\
\textbf{Kurs} \> TINF17IN\\
\textbf{Betreuer der Hochschule} \> Alfred Becker\\
\textbf{Gutachter der Hochschule} \> Alfred Becker\\
\end{tabbing}
\end{spacing}
%Seitennummerierung ausschalten
\pagenumbering{gobble}
\newpage

\section*{Selbstständigkeitserklärung}

\bigskip

Ich versichere hiermit, dass ich meine Bachelorarbeit (bzw. Studien- und Projektarbeit) mit dem Thema:

\smallskip

%% eigentlich hier \titel verwenden statt duplicated titel, aber umbruch erzwingt
%% doppeltes ausschreiben des titels
\texttt{Entwicklung eines Chatsystems auf Basis von XMPP}

\smallskip

selbstständig verfasst und keine anderen als die angegebenen Quellen und Hilfsmittel benutzt habe.

\bigskip

Ich versichere zudem, dass die eingereichte elektronische Fassung mit der gedruckten Fassung übereinstimmt.*

\bigskip

\begin{small}

* falls beide Fassungen gefordert sind

\bigskip

\bigskip

\noindent\begin{tabular}{ll}
\makebox[2.5in]{\hrulefill} & \makebox[2.5in]{\hrulefill}\\
Ort, Datum & Unterschrift
\end{tabular}
\end{small}

\newpage

%abstract text
\section*{Abstract}

\newpage

%inhaltsverzeichnis
	% Inhaltsverzeichnis
	\cleardoublepage
	\begin{spacing}{1.1}
		\begingroup
		
			% auskommentieren für Seitenzahlen unter Inhaltsverzeichnis
			\renewcommand*{\chapterpagestyle}{empty}
			\pagestyle{empty}
			
			
			%\setcounter{tocdepth}{1}
			%für die Anzeige von Unterkapiteln im Inhaltsverzeichnis
			\setcounter{tocdepth}{2}
			
			\tableofcontents
			\clearpage
		\endgroup
	\end{spacing}

%% new header/footer settings
\renewcommand{\sectionmark}[1]{\markright{\thesection\ #1}} % make header rightmark
\fancypagestyle{fancyheadlines}{
\pagenumbering{arabic}
\fancyhf{}
\lhead{\slshape\rightmark}
%%\rhead{\slshape\nouppercase{\leftmark}}
\renewcommand{\headrulewidth}{0.4pt}
%\lfoot{\slshape DHBW Stuttgart | Lukas Priester, Oliver Klapper}
\cfoot{\thepage}
\renewcommand{\footrulewidth}{0.4pt}
}

% Redefine the plain page style, show only page number in figure,table,...contents
% and chapter pages
\fancypagestyle{plain}{%
  \fancyhf{}%
  %\lfoot{\slshape DHBW Stuttgart | Lukas Priester, Oliver Klapper}%
  \cfoot{\thepage}
  \renewcommand{\headrulewidth}{0pt}% Line at the header invisible
  \renewcommand{\footrulewidth}{0.4pt}% Line at the footer visible
}


\newpage
\pagenumbering{Roman}


%abkürzungsverzeichnis
\cleardoublepage
\addcontentsline{toc}{chapter}{Abkürzungsverzeichnis}
\chapter*{Abkürzungsverzeichnis}
\begin{acronym}[YTMMM]
\setlength{\itemsep}{-\parsep}

\acro{XMPP} {Extensible Messaging and Presence Protocol}
\acro{NLP} {Natural Language Processing}
\acro{MUC} {Multi User Chat}
\acro{API} {Application Programming Interface}
\acro{HTML}{HyperText Markup Language}
\end{acronym}

%abbildungsverzeichnis
\cleardoublepage
\addcontentsline{toc}{chapter}{\listfigurename}
\listoffigures
\newpage
%tabellenverzeichnis
\cleardoublepage
\addcontentsline{toc}{chapter}{\listtablename}
\listoftables
\newpage
%listingverzeichnis
\cleardoublepage
\addcontentsline{toc}{chapter}{\lstlistingname}
\lstlistoflistings
\newpage

\begin{onehalfspacing}

%% header and footer settings
\pagestyle{fancyheadlines}

\chapter{Einleitung}
\label{chap:Einleitung}

Social Media (deutsche Übersetzung: \glqq soziale Medien\grqq{}) ist ein aktuelles und wichtiges gesellschaftliches Thema. Die vielfältigen und breitgefächerten Nutzungsmöglichkeiten beeinflussen das Privat- und das Berufsleben \cite{gabriel2017social}. Gabriel und Röhrs definieren Social Media in \cite{gabriel2017social} als die Verwendung digitaler Medien unter Einsatz computergeschützter Technologien, das heißt, von Hardware- und Softwaresystemen. Sie definieren den Nutzen von Social Media darin, dass Menschen Informationen suchen, erstellen, verteilen und austauschen können. Es gibt eine große Anzahl an unterschiedlichen Definitionen von Social Media. Nach Liu Yinyuan ist Social Media längst wichtiger Bestandteil des Unternehmensmarketings in Deutschland. In seinem Werk \glqq Social Media in China\grqq{} \cite{liu2016social} beschreibt er, dass in Unternehmen nicht mehr über die grundsätzliche Frage debattiert wird, ob Social Media für das Unternehmensmarketing eingesetzt werden soll, sondern wo und wie der Einsatz zielgerichtet erfolgen kann. Heutige Unternehmen sind gekennzeichnet von computergestützten Anwendungssystemen, die in allen Funktionsbereichen zur Planung, Steuerung und Kontrolle der Geschäftsprozesse und zu ihrer Verwaltung eingesetzt werden. Sie setzen dieses System in der B2B-Kommunikation (Business-To-Business-Kommunikation) ein. Immer wichtiger werden \textbf{Instant Messaging Systeme}, wie zum Beispiel Whatsapp, Telegram, iMessage und Jabber, die dazu dienen intern Informationen schnell im Unternehmen zu verbreiten, hoch verfügbar zu machen und um Geschäftsprozesse standortunabhängig steuern zu können \cite{gabriel2017social}. Zusätzlich werden \textbf{Instant Messaging Systeme} auch zur externen Kommunikation benutzt. Nach Gabriel und Röhrs ist es möglich, dass mehrere Unternehmen mit Hilfe von innovativen Kommunikationstechniken zur Erreichung eines gemeinsamen Ziels besser kooperieren können. Außerdem ist die schnelle und direkte Kontaktaufnahme von Kunden über ein Messaging System zum Support eines Unternehmens eine einfache und schnelle Möglichkeit, Fragen zum Produkt ohne langes Warten in der Hotline zu stellen. Nach \cite{b2bmehner} werden Nachrichten von \textbf{Instant Message Systemen} im Gegensatz zu einer E-Mail in Echtzeit übertragen und dem Empfänger direkt zugestellt. Laut einer Statistik von \textbf{statista} benutzen 1,5 Milliarden Nutzer in Deutschland pro WhatsApp pro Monat \cite{statistaIMS}. Der Anteil der Nutzer von Whatsapp in Deutschland beträgt 75 Prozent an der Gesamtbevölkerung \cite{statistaIMS}. Die Statistik zeigt, dass \textbf{Instant Messaging Systeme} eine Möglichkeit in der Zukunft darstellen Kunden direkter anzusprechen oder Support zu gewährleisten. Zusätzlich zu Nachrichtendiensten werden künstliche Intelligenzen und Algorithmen benötigt, die Nutzer unterstützen oder gesammelte Daten von Benutzern charakterisieren oder interpretieren können. Algorithmen werden verwendet, um zum Beispiel die emotionale Befindlichkeit anhand eines Text einzuordnen, um Suizid-Gedanken frühzeitig zu erkennen \cite{stasytisIME}. Ein weiteres Anwendungsszenario sind Text-to-Speech (deutsche Übersetzung: Text-zu-Sprache) Funktionalitäten, bei denen gesprochene Worte des Nutzers durch Algorithmen in Text umgewandelt werden, sodass ein Nutzer die Nachricht nicht mehr aktiv eintippen muss.

\section{Aufgabenstellung}
\label{sec:Aufgabenstellung}

Die konkrete Aufgabe ist es, einen Chatserver auf Basis des \textbf{\ac{XMPP}}-Protokolls in Betrieb zu nehmen über den sich mehrere Chat-Clients authentifizieren und verschlüsselt Nachrichten austauschen können. Durch eine gründliche Recherche soll eruiert werden, welche Chatserver sich hierfür eignet und warum dieser Chatserver für das Projekt verwendet wird. Der Nachrichtenaustausch soll Ende-zu-Ende verschlüsselt erfolgen. Zudem sollen auch Gruppenchats möglich sein. Eine Zustellung der Nachrichten in Echtzeit soll implementiert werden. Benutzer sollen Nachrichten über eine Web-Oberfläche eingeben und empfangen können. Es soll ein Prototyp eines Machine Learning Algorithmus der Kategorie \textbf{\ac{NLP}} implementiert werden und in die Web-Oberlfäche integriert werden.

\section{Ziele der Arbeit}
\label{sec:Ziele}

Die Ziele der Arbeit sind es, eine geeignete Netzwerkumgebung und einen lauffähigen Chatserver auf Basis von \ac{XMPP} zu implementieren. Es sollen Kenntnisse und Erfahrungen mit dem Umgang des \ac{XMPP} Protokolls gesammelt werden und sich mit dem Aufbau des Nachrichtenprotokolls auseinandergesetzt werden. Es soll eine Web-Oberfläche implementiert werden, die einen \ac{MUC} unterstützt und über den Chatserver Nachrichten versendet und empfängt. Durch intensiver Recherche soll ein geeignetes Framework gefunden werden, mit dem eine Web-Oberfläche zusammen mit einer \ac{API} programmiert werden kann.


\section{Stand der Technik}
\label{sec:StandDerTechnik}

\newpage

\chapter{Datenschutzrechtliche Aspekte}
\label{chap:Datenschutz}

\newpage

\chapter{Anforderungen}
\label{Anforderungen}

\newpage

\chapter{Theoretische Grundlagen}
\label{chap:Theorie}

\section{XMPP-Protokoll}
\label{sec:XMPP}

\section{Python}
\label{sec:Python}

\section{Chatserver}
\label{sec:Chatserver}

\begin{lstlisting}[language=Python,caption=Example Listing Python,label={lst:Example}]
"""The first step is to create an SMTP object, each object is used for connection 
with one server."""

import smtplib
server = smtplib.SMTP('smtp.gmail.com', 587)

#Next, log in to the server
server.login("youremailusername", "password")

#Send the mail
msg = "
Hello!" # The /n separates the message from the headers
server.sendmail("you@gmail.com", "target@example.com", msg)
\end{lstlisting}


\section{Datenbank}
\label{sec:Datenbank}

\newpage

\chapter{Umsetzung und Implementierung}
\label{chap:Umsetzung}

\section{Vorbereitung}
\label{sec:Vorbereitung}

\section{Implementierung des Chatservers ejabberd}
\label{sec:ChatserverEntwicklung}

\subsection{Anbindung des Chatservers}
\label{subsec:Anbindung}

\subsection{Security}
\label{subsec:Security}

\subsection{Konfiguration von ejabberd}
\label{subsec:Konfiguration}

\newpage

\section{Entwicklung des Chat-Clients}
\label{sec:ClientEntwicklung}

\ac{HTML}

\subsection{Entwicklung der GUI}
\label{subsec:EntwicklungGUI}

\subsection{Implemintierung des Backends}
\label{subsec:Backend}

\newpage

\chapter{Datenanalyse}
\label{chap:Datenanalyse}

\newpage

\chapter{Ergebnis}
\label{chap:Ergebnis}

\newpage

\chapter{Fazit}
\label{chap:Fazit}


\end{onehalfspacing}
\newpage

\bibliography{Literatur}
\newpage
%setze Anhang mit appendix
%addpart, sodass Anhang im Inhaltsverzeichnis ohne Buchstaben auftaucht
\appendix
\addpart{Anhang}

\end{document}